% ---------------------------------------------------------------------------
%  UNIT II -- Chapters 8 and 9
% ---------------------------------------------------------------------------
\chapter{Unit II --- Harmony in the Family and in the Society: From Family
  Order to World Family Order}
\chaptermark{Unit II --- Chapters 8 and 9}

\section{A. Relationship, Justice and the State Today (Chapter 8)}

% ------------------------------------------------------------------------ 27
\Q{27}{``Family is a natural laboratory to understand human relationships''
  --- elaborate. \emph{(Also set as:} Relationship IS, and it exists between
  one jeevan and the other jeevan. Examine.\emph{)}}{\tg{CIE 2 QB}%
  \tg{CIE-II 2025}\mmk{5}}

\ans The family is the \textbf{basic unit of human interaction}. Each one of
us is naturally a part of a family that includes father, mother, brothers and
sisters, and then grandparents, aunts, uncles, cousins. These relationships
are a reality of our life --- we are \emph{born into} them; we have not
created them.

Here is the set of proposals the text asks us to verify:

\begin{apts}
\item \textbf{Relationship IS, and it exists between the Self (`I') and the
      other Self (`I').} It is not possible to create the relationships that
      are existent in a family; we are naturally born into them. It is the
      person's self which is primarily related to the other person's self. The
      Body is only a means to express or receive our relationship --- the Body
      is incapable of understanding as well as having feelings. For example, a
      mother feels related to the child she has given birth to; the body of
      the child has its source in the body of the mother, but neither body has
      feelings --- it is the Self of the mother and the child who feel
      connected.
\item \textbf{The Self (`I') has feelings in a relationship, and these
      feelings are between `I' and `I'.} Who wants trust in relationship ---
      you or the Body? The answer is, I want trust; and from the other `I'.
      There is no part of the body that wants trust or respect.
\item \textbf{These feelings in the Self are definite --- they can be
      identified with definiteness.} They are the values characterising
      relationships: trust, respect, affection, and so on.
\item \textbf{Recognising and fulfilling these feelings leads to mutual
      happiness in relationship.}
\end{apts}

\textbf{Why the family is a laboratory.} By living in relationships at all
times in the family, we get the occasion to gain the assurance that the other
person is an \textbf{aid to me and not a hindrance}. In the family we learn to
recognise relationship, recognise the definite feelings or values, and learn
how to fulfil them. The evaluation that takes place mutually in these close
relationships, leading to mutual happiness, instils a confidence in us that we
can live the right way with human beings. The family is thus a laboratory of
sorts in which we live our understanding. On getting assured, it becomes easy
to see that \textbf{society is an extension of family}, and that it is
possible to live in harmony with every human being --- thus laying the
foundation for an undivided human race, from \emph{family order} to
\emph{world family order}.

% ------------------------------------------------------------------------ 28
\Q{28}{What is `justice'? What are its four elements? Is it a continuous or a
  temporary need?}{\tg{SEE}\tg{CIE 2 QB}\tg{Scheme}\mmk{5}}

\ans \textbf{Justice is the recognition of values (the definite feelings) in
relationship, their fulfilment, and the right evaluation of the fulfilment
resulting in mutual happiness (Ubhay-tripti).} Justice concerns itself with
the proper ordering of things and people within a society.

\begin{keybox}
\textbf{The four elements of justice:}
\begin{enumerate}[label=\textbf{\color{ocre}\arabic*.},leftmargin=2.0em,
  itemsep=1pt,topsep=3pt,parsep=0pt]
\item \textbf{Recognition of values} --- recognising the definite feelings in
      the relationship
\item \textbf{Fulfilment} of those values
\item \textbf{Evaluation} of the fulfilment
\item \textbf{Mutual happiness} ensured
\end{enumerate}
\vspace{3pt}
When all four are ensured, justice is ensured. \textbf{Mutual fulfilment is
the hallmark of justice.}
\end{keybox}

\begin{center}\small\sffamily\color{slate}
Self (`I') ---relationship--- Self (`I') \ $\rightarrow$ \ Recognise the
VALUES \ $\rightarrow$ \ Fulfilment of VALUES \ $\rightarrow$ \ Evaluate the
fulfilment \ $\rightarrow$ \ Mutual happiness
\end{center}

\textbf{Is it continuous or temporary?} Justice is a \textbf{continuous need}.
Justice is essential in all relationships --- be it with the small kid in your
house, your old grandpa, the maid in the house, your fast friends or your
distant relations. We want justice not on a few occasions but \emph{every
moment}, and in every relationship, not just a few. We need to grow up in
relationships to ensure the continuity of justice in all our relationships. It
is also worth asking whether justice will get ensured in the family or in
courts of law --- in justice there is \emph{mutual fulfilment} for both
parties, which a judgement and punishment cannot deliver.

% ------------------------------------------------------------------------ 29
\Q{29}{What is the meaning of justice in human relationships? How does it
  follow from family to world family?}{\tg{CIE 2}\tg{SEE}\tg{Scheme}\mmk{5}}

\ans Justice is the recognition of values (the definite feelings in a
relationship), their fulfilment, and the right evaluation of the fulfilment
resulting in mutual happiness. There are four elements: recognition of values,
fulfilment, evaluation, and mutual happiness ensured. When all four are
ensured, justice is ensured; mutual fulfilment is the hallmark of justice.

\textbf{From family to world family.} Justice \textbf{starts from the family
and slowly expands to the world family}. The child gets an understanding of
justice in the family. With this understanding, he\,/\,she goes out into
society and interacts with people. All of us are children at some point of
time and grow into adults. In the family we learn to recognise relationship,
recognise the definite feelings or values, and learn how to fulfil them. The
evaluation that takes place mutually in close relationships, leading to mutual
happiness, instils a confidence in us that we can live the right way with
human beings; unless this confidence is ensured, we remain shaky in
relationships.

\textbf{If the understanding of justice is ensured in the family, there will
be justice in all the interactions we have in the world at large.} If we do
not understand the values in relationships, we are governed by our petty
prejudices and conditionings. We may treat people as high or low based on
their body (particular caste, or sex or race or tribe), on the basis of the
wealth one possesses, or the belief systems that one follows. All this is a
source of injustice and leads to a \emph{fragmented society}, while our
natural acceptance is for an \emph{undivided society and universal human
order}.

Having explored the harmony in human beings, we are able to explore the
harmony in the family. This enables us to understand the harmony at the level
of society and nature\,/\,existence. And this is the way the harmony in our
living grows --- we slowly get the competence to live in harmony with all
human beings.

% ------------------------------------------------------------------------ 30
\Q{30}{What is the outcome when we try to identify relationships based on the
  exchange of physical facilities?}{\mmk{5}}

\ans We are unable to see ourselves as the co-existence of the Self (`I') and
the Body. As a result we see ourselves as a body and the other as a body, and
we subsequently \textbf{reduce our relationships and the feelings in the
relationship to the level of the body}. We tend to assume that we have
relationship with our blood-related family members only.

As a result of this mistaken assumption, we have \textbf{reduced our
expectations in relationships to the mere fulfilment of physical facilities}.
We evaluate all our relationships in terms of material things like money,
property, etc. In short, the purpose of relationship has been reduced to
physical and material needs and their exchange. Hence we feel that working for
physical facilities alone is enough, or we assume that as long as we
accumulate physical facilities and provide the same to the other, the
relationship is automatically fulfilled.

\textbf{The outcome.} Suppose your father earns enough money and ensures that
your physical needs are taken care of, but does not spend time with you, does
not care for you, or instead behaves badly with you --- would you feel
satisfied? The answer is NO. Nowadays we also hear of youngsters earning a lot
of money who, instead of taking care of their parents and fulfilling their
needs of feelings at the level of `I' (trust, respect, affection), just put
their parents in some old-age home. The parents have plenty to eat, good
clothes, a big TV, servants --- but this is not fulfilling for them, since the
needs of the `I', the feelings in `I', have been totally ignored.

\textbf{The fact is:} what we need first is the \textbf{right understanding},
and this is not ensured by having money. Secondly, we need the \textbf{feelings
in relationships} to be fulfilled, which is also not ensured by having money.

\section{B. The Values in Human Relationships (Chapter 8)}

% ------------------------------------------------------------------------ 31
\Q{31}{Enumerate the important values which lie at the base of good
  relationships. \emph{(Also set as:} List the foundation value and the
  complete value in human relationship, and explain each with one
  example.\emph{)}}{\tg{CIE 2 QB}\mmk{5}}

\ans Being in relationship, we have feelings for the other. These feelings
cannot be replaced by any material or physical things. They are definite and
they are the values in a relationship. There are \textbf{nine} such values:

\renewcommand{\arraystretch}{1.3}
\arrayrulecolor{rulegrey}
\begin{longtable}{C{7mm} P{45mm} P{95mm}}
\rowcolor{slate}\hdr{\#} & \hdr{Value} & \hdr{Meaning (as defined in the text)}
  \\ \hline \endfirsthead
\rowcolor{slate}\hdr{\#} & \hdr{Value} & \hdr{Meaning (as defined in the text)}
  \\ \hline \endhead
\qn{1} & \textbf{Trust} (Vishvas) & To be assured that each human being
  inherently wants oneself and the other to be happy and prosperous.
  \textbf{The foundation value} (adhara mulya). \\ \hline
\qn{2} & \textbf{Respect} (Sammana) & Right evaluation --- to be evaluated as
  I am, without over-, under- or otherwise evaluation. \\ \hline
\qn{3} & \textbf{Affection} (Sneha) & The feeling of being related to the
  other; the acceptance of the other as one's relative. \\ \hline
\qn{4} & \textbf{Care} (Mamata) & The feeling to nurture and protect the body
  of our relative. \\ \hline
\qn{5} & \textbf{Guidance} (Vatsalya) & The feeling of ensuring right
  understanding and right feelings in the other (my relative). \\ \hline
\qn{6} & \textbf{Reverence} (Shraddha) & The feeling of acceptance of
  excellence in the other. \\ \hline
\qn{7} & \textbf{Glory} (Gaurava) & The feeling for someone who has made
  efforts for excellence. \\ \hline
\qn{8} & \textbf{Gratitude} (Kritagyata) & The feeling of acceptance for those
  who have made effort for my excellence. \\ \hline
\qn{9} & \textbf{Love} (Prema) & The feeling of being related to all.
  \textbf{The complete value} (purna mulya). \\ \hline
\end{longtable}

\textbf{Foundation value --- Trust}, with example: In a family, when a member
fails to do something for me, if I am assured that at the level of intention
he\,/\,she means well for me and only the competence was lacking, I do not get
hurt; I become a help to him\,/\,her. A relationship without trust results in
opposition; the relationship itself gets shaken up. Lack of trust is what
ultimately leads to extreme situations like war.

\textbf{Complete value --- Love}, with example: It starts with identifying
that I am related to one other human being (the feeling of affection) and
slowly expands to the feeling of being related to all human beings. When I am
able to feel related even to a stranger on the road --- recognising that at
the level of `I' he\,/\,she is exactly like me --- that is love. This feeling
of love lays down the basis of an \textbf{Undivided Society}.

Two useful points the text stresses: the basic values to be understood are
\textbf{trust and respect} --- if we have these, the remaining values flow
quite naturally. And \textbf{only Care (mamata) requires physical facilities};
for all the other feelings, what we essentially need is their proper
understanding.

% ------------------------------------------------------------------------ 32
\Q{32}{What do you understand by trust? Differentiate between intention and
  competence with examples. Why is this distinction important?}%
  {\tg{CIE 2}\tg{SEE}\tg{CIE-II 2025}\mmk{5}}

\ans \textbf{Trust (vishvas)} is defined as: ``\emph{To be assured that each
human being inherently wants oneself and the other to be happy and
prosperous.}'' We feel assured of the other person when we are sure that the
other wants to work for my happiness and prosperity. Whenever I feel the other
will deny my happiness and\,/\,or prosperity, I am afraid of that person.
Trust is the \textbf{foundation value} (adhara mulya) of relationship.

To understand it, examine four proposals: (1) I want to make myself happy;
(2) I want to make the other happy; (3) The other wants to be happy; (4) The
other wants to make me happy. Answering the fourth with deep exploration is
the basis of gaining trust in the other. On examining, we find there are
\textbf{two parts} in this exploration:

\begin{cmptable}{Intention (natural acceptance) --- \emph{wanting to}}%
                {Competence (ability to fulfil) --- \emph{being able to do}}
1a) I want to be happy \ \checkmark\newline
2a) I want to make the other happy \ \checkmark\newline
3a) The other wants to be happy \ \checkmark\newline
4a) The other wants to make me happy \ \checkmark \ (after exploration) &
1b) I am always happy \ ?\newline
2b) I always make the other happy \ ?\newline
3b) The other is always happy \ ?\newline
4b) The other always makes me happy \ ? \\ \hline
\textbf{What we really want to be} & \textbf{What we are} \\ \hline
\end{cmptable}

\textbf{The mistake we make.} When we judge ourselves, we judge on the basis
of our \emph{intention}; when we judge the other, we judge on the basis of
his\,/\,her \emph{competence}. We say ``I wanted to do well, but I could
not'' --- but for the other we say ``He did not want to do well''. ``Wanting
to'' is the intention; ``could not'' is the lack of competence.

\textbf{Example.} You are walking on your college campus and your close friend
walks by from the other direction. You look at him and smile, but he barely
notices you and keeps walking with his head down. You feel angry and
disappointed, and assume he wants to ignore you. Later you find out that he
was disturbed since he had lost his wallet. You immediately feel alright and
are not angry any more. What happened? You \textbf{doubted your friend's
intention}. It was not that he intended to ignore you; only that he was
preoccupied. Once you knew this, you realised it was not his intention but
only \textbf{his competence that was lacking} at that moment.

\textbf{Why the distinction is important.} Intention is \textbf{always
correct} and is the same in all of us; it is only the competence that is
lacking, which can be improved by right understanding. When we are assured of
the intention of the other and find that the competence is lacking, we
\textbf{become a help} to the other. When we doubt the intention of the other,
we \textbf{get into opposition}. You cannot have a problem in relationship
unless you have ended up doubting the intention of the other person, no matter
how close you are to them.

% ------------------------------------------------------------------------ 33
\Q{33}{In our behaviour, we generally observe our intention and others' lack
  of competence. Does it lead to mutual happiness? What is the alternative?
  Explain with an example.}{\tg{CIE 2 QB}\tg{Scheme}\mmk{5}}

\ans \textbf{No, it does not lead to mutual happiness.} We trust our own
intention while we are not ready to trust the other's intention. It is the
same for the others as well --- while the other trusts his\,/\,her own
intentions, he\,/\,she does not trust mine. Hence \textbf{mistrust is born and
we deny the relationship}.

When we are judging ourselves we are judging on the basis of our intention,
whereas when we are judging the other we are judging him on the basis of his
competence. We are sure that we want to make the other happy, but we are not
sure that the other wants to make us happy. We find that while we look at our
intention we are sure of it; we are not sure of the other's intention. We are
actually seeing their competence and making a conclusion on their intention
--- we say ``I wanted to do well, but I could not'', but for the other we say
``He did not want to do well''.

\textbf{The alternative.} We can see that we are not able to fulfil our
intentions in terms of our competence at all times; it is the same for the
other as well. We want to be related to the other, and we want the other to be
related to us, irrespective of who this other is. If we have trust in the
other, we are able to see the other as a \textbf{relative and not as an
adversary}. We then become ready to \textbf{become a help to the other}.
Intentions are always correct; it is only the competence that is lacking,
which can be improved by the right understanding.

\textbf{Example.} Suppose a friend promises to bring your notes to class and
forgets. If I doubt his intention (``he did not want to help me''), I get
irritated, we argue, and both of us are unhappy --- there is no mutual
happiness. If instead I am assured of his intention and see that only his
competence (memory, planning) was lacking, I do not get hurt; I remind him and
help him organise better next time. He is comfortable and so am I --- this is
mutual happiness, and his competence also improves.

% ------------------------------------------------------------------------ 34
\Q{34}{There is a common saying: ``If you trust everybody, people will take
  undue advantage of you.'' What is the basic error in this statement?
  \emph{(Also set as:} How can I trust a stranger?\emph{)}}%
  {\tg{CIE 2 QB}\tg{Scheme}\mmk{5}}

\ans \textbf{The basic error is that if we trust everybody, people will
\emph{not} take undue advantage of us.} On the contrary, it gives us
\textbf{inner strength} and we become far more effective in interacting with
and dealing with different people. This is simply because we already are
sitting with the knowledge of what the person truly wants, truly intends, even
though the person may not know this himself\,/\,herself.

Hence our ability to interact with people becomes far more effective, and in
the process \textbf{we don't get hurt, we don't get disturbed, and we end up
becoming an aid to the other}. In other words, becoming aware, having the
right understanding, and living with the assurance in the relationship does
not mean becoming ``stupid'' --- it only makes us \textbf{more competent}.
People can take advantage of you only if you \emph{do not} have the right
understanding, which is the state we are in today.

Further, what is being said is that we have \textbf{trust in the intention} of
everyone; but when it comes to \textbf{making a programme} with someone, I
evaluate my competence, I evaluate his competence, and make the programme
accordingly. This makes me more effective than if I did it otherwise, i.e. by
doubting his intention.

\textbf{``How can I trust a stranger?''} If you are able to see the
relationship with the person at the level of `I', you will see that the other
person is also like you. He\,/\,she too has natural acceptance for the same
things as you, wants to make himself\,/\,herself happy, and wants to make you
happy at the level of intention --- just as you do. But he\,/\,she may be
unaware of this fact, just as you were, and hence may be interacting with you
on the basis of your competence. \textbf{If we interact with or evaluate the
other person at the level of competence only, then there cannot be continuity
of trust.} In that case we end up doubting the other person, which causes a
sense of opposition in us; and since opposition is not naturally acceptable,
it creates a contradiction in us. Hence the way out is to relate to the other
person, to be able to see that at the level of natural acceptance we are the
same. We can then interact with the person based on their competence, and also
help them improve their competence.

\begin{keybox}
So, you can trust anyone \emph{for the intention part}. But don't assume that
his\,/\,her desires, thoughts and expectations are going to be right ---
he\,/\,she may lack competence. \textbf{The competence is to be evaluated
before you make a programme with the other.}
\end{keybox}

% ------------------------------------------------------------------------ 35
\Q{35}{What is the meaning of respect? What is the basis of respect for a
  human being? Do you see that the other human being is also similar to
  you?}{\tg{CIE 2 QB}\tg{SEE 2025}\mmk{5}}

\ans \textbf{Respect (sammana) means ``Right Evaluation''} --- to be evaluated
as I am. Usually, however, we make mistakes in evaluation in three ways:

\renewcommand{\arraystretch}{1.3}
\arrayrulecolor{rulegrey}
\begin{longtable}{P{45mm} P{40mm} P{62mm}}
\rowcolor{slate}\hdr{Mistake} & \hdr{Meaning} & \hdr{Example} \\ \hline
\endfirsthead
\rowcolor{slate}\hdr{Mistake} & \hdr{Meaning} & \hdr{Example} \\ \hline
\endhead
\textbf{Over evaluation} (adhi-mulyana) & To evaluate more than what it is &
  Your father tells guests, ``My son is the greatest scholar in India!'' ---
  you feel uncomfortable when wrongly flattered. \\ \hline
\textbf{Under evaluation} (ava-mulyana) & To evaluate less than what it is &
  ``My son is a good for nothing; he must be the laziest person in all of
  India!'' --- you feel uncomfortable when condemned. \\ \hline
\textbf{Otherwise evaluation} (a-mulyana) & To evaluate otherwise than what it
  is & ``You donkey! Can't you even understand this much?'' --- you are
  evaluated as something else, since you are a human being and not something
  else. \\ \hline
\end{longtable}

Any kind of over, under or otherwise evaluation makes us uncomfortable; we
feel `disrespected'. \textbf{We say we have been disrespected when we are
wrongly evaluated.}

\textbf{The basis for respect.} A human being is a co-existence of Self (`I')
and Body. Right evaluation is on the basis of the acceptance of this
co-existence. If we respect a human being on the basis of `I', the following
are true for every human being:

\begin{apts}
\item I want continuous happiness and prosperity --- \emph{the other too wants
      to be continuously happy and prosperous}.
\item To be happy, I need to understand and live in harmony at all four levels
      of my living --- \emph{the other also needs to understand and live in
      harmony at all four levels of his\,/\,her living}.
\item The activities in me (`I') are continuous --- we can check this for
      desire, thought and expectation --- \emph{it is the same for the other
      `I' as well}.
\end{apts}

Based on these three evaluations we conclude:

\begin{abul}
\item Our \textbf{basic aspiration} is the same --- we both want continuous
      happiness and prosperity.
\item Our \textbf{programme of action} is the same --- to understand and live
      in harmony at all four levels of our living.
\item Our \textbf{potential} is the same --- the activities and powers of the
      self are continuous and the same in both of us at the level of `I'.
\end{abul}

\begin{keybox}
$\Rightarrow$ \ \textbf{The other is similar to me.}
\end{keybox}

Yes --- when we are able to see that the other is similar to me, we are able
to recognise the feeling of respect in the relationship. If not, we either
hold ourselves more or less than the other, and this only leads to
differentiation. The difference between me and the other can only be at the
level of \emph{understanding} (not information), and this manifests as a
meaningful responsibility, not as a criterion for superiority or inferiority:
if the other has better understanding than me, I want to understand from the
other; if the other has less understanding than me, I accept the
responsibility to improve the understanding of the other.

% ------------------------------------------------------------------------ 36
\Q{36}{How do we differentiate in relationships on the basis of body, physical
  facilities or beliefs? What problems do we face because of such
  differentiation? \emph{(Also set as:} ``Discrimination leads to acrimony in
  relationships.'' Explain.\emph{)}}%
  {\tg{CIE 2}\tg{SEE}\tg{CIE-II 2025}\tg{Scheme}\mmk{5}}

\ans Instead of respect being a basis of similarity or of \emph{right
evaluation}, we have made it into something on the basis of which we
differentiate. We differentiate on three bases:

\renewcommand{\arraystretch}{1.3}
\arrayrulecolor{rulegrey}
\begin{longtable}{P{22mm} P{60mm} P{65mm}}
\rowcolor{slate}\hdr{Basis} & \hdr{How we differentiate} &
  \hdr{Problems created (as per the scheme)} \\ \hline \endfirsthead
\rowcolor{slate}\hdr{Basis} & \hdr{How we differentiate} &
  \hdr{Problems created (as per the scheme)} \\ \hline \endhead
\textbf{Body} & \textbf{Sex\,/\,Gender} --- respect males more than females
  (or the reverse), ignoring that being male or female is an attribute of the
  body, not of `I'. &
  Issue of women's rights; women protesting and demanding equality in
  education, in jobs and in people's representation. People are insecure and
  afraid of one another based on their gender. \\ \hline
 & \textbf{Race} --- skin colour, Mongolian\,/\,Aryan\,/\,Dravidian race, or
  caste taken as high and low. &
  Many movements and protests against racial discrimination and demands for
  equality; racial attacks; movements against caste discrimination; people
  living in fear of racism, racist attacks, casteism and discrimination.
  \\ \hline
 & \textbf{Age} --- notions such as ``one must respect elders''; age relates
  to the body, not to `I'. &
  Protests and movements demanding equal rights for children on the one hand
  and rights for elderly people on the other; the \emph{generation gap}, a
  source of tension in many families. \\ \hline
 & \textbf{Physical strength} --- treating the stronger person differently. We
  think we are respecting, while it is actually fear. &
  Many programmes, awards and titles based on physical strength, which have
  nothing to do with how the person is at the level of `I'. \\ \hline
\textbf{Physical facilities} & \textbf{Wealth} --- the ``rich person'' gets
  idolised; we never check whether such people feel prosperous or are happy. &
  Class struggle and movements to do away with class-differentiation; many
  people suffering from a lack of self-esteem, and some even committing
  suicide. \\ \hline
 & \textbf{Post} --- respect on the basis of a person's position; the post
  wrongly evaluated as the mark of excellence. &
  Protests against high-handed government officials. At the level of the
  individual, it leads to depression, since those who do not qualify for a
  post feel they will not get respect in society. \\ \hline
\textbf{Beliefs} & \textbf{`Isms'} --- capitalism, socialism, communism, etc.;
  belief systems adopted since childhood. &
  Fights, turmoil, terrorism and war; people converting from one `ism' to
  another in order to get more respect. \\ \hline
 & \textbf{Sects} --- a set of beliefs reflected in traditions and practices;
  only those with a similar belief system are considered `own' and worthy of
  respect. &
  Countless religions and sects, each with its own movement to ensure there is
  no discrimination against people of their belief; demands for special
  provisions in jobs and in education; clashes that even turn violent.
  \\ \hline
\end{longtable}

\textbf{The way out.} To move beyond differentiation we have to understand the
human being as the co-existence of Self (`I') and Body, and then base our
evaluation on the Self (`I'), where we find we are similar to the other in
natural acceptance, programme of action and potential. This becomes the basis
of the feeling of respect. We differ from the other only in terms of
competence, and there we either learn from the other or take the
responsibility of helping the other improve.

% ------------------------------------------------------------------------ 37
\Q{37}{What is the difference between respect and disrespect? Which of the two
  is naturally acceptable to you? \emph{(Also set as:} List the differences
  between respect and differentiation.\emph{)}}{\tg{CIE 2}\tg{CIE 2 QB}%
  \tg{SEE}\mmk{5}}

\ans
\begin{cmptable}{Respect (Right Evaluation)}{Disrespect \,/\, Differentiation}
Means \textbf{right evaluation} --- to evaluate the other as he\,/\,she is. &
Means \textbf{wrong evaluation} --- over evaluation, under evaluation, or
  otherwise evaluation. \\ \hline
Evaluation is done at the level of the \textbf{Self (`I')}. &
Evaluation is done at the level of the \textbf{Body, physical facilities or
  beliefs}. \\ \hline
Based on \textbf{similarity} --- the other is similar to me in aspiration,
  programme of action and potential. &
Based on \textbf{difference} --- sex, race, age, physical strength, wealth,
  post, `isms', sects. \\ \hline
Leads to a feeling of relationship, mutual fulfilment and ease in both. &
Leads to a feeling of opposition, discomfort, and problems in relationship.
  \\ \hline
It is \textbf{continuous} --- once we have this feeling out of understanding,
  it remains there all the time; it is the state of harmony between one human
  and the other. &
It is \textbf{unstable} and changes with the person's body, possessions or
  position, and hence needs constant re-establishment. \\ \hline
Does not require the other to do anything special; the fact that you are human
  is enough for me to respect you. &
Makes everyone run around trying to become ``special'' to earn respect ---
  which fetches only attention, not respect. \\ \hline
Its outcome is an \textbf{undivided society}. &
Its outcome is a \textbf{fragmented society}. \\ \hline
\end{cmptable}

\textbf{Which is naturally acceptable? Respect.} Verify it for yourself: what
is naturally acceptable to you --- the feeling of respect or disrespect for
yourself? And for the other? You will find that each one of us has an
acceptance for the feeling of respect. Just as we desire this, the other also
expects the same. \textbf{Every human being wants to respect and be
respected.} No human being finds disrespect acceptable in relationship.

% ------------------------------------------------------------------------ 38
\Q{38}{Define affection. How does affection lead to harmony in the family?
  What is the role of physical facilities in the fulfilment of this
  feeling?}{\tg{CIE 2}\mmk{5}}

\ans \textbf{Affection (sneha) is the feeling of being related to the other.}
Affection comes when I recognise that we both want to make each other happy
and that both of us are similar. Then, for the first time, I feel that I am
related to the other, that the other is a relative of mine. ``\emph{This
feeling of acceptance of the other as one's relative is the feeling of
affection or sneha in relationship.}''

\textbf{How it leads to harmony in the family.}

\begin{abul}
\item This feeling comes only if \textbf{Trust and Respect are already
      ensured}. Without trust and respect you feel the other is trying to make
      you unhappy and does not wish well for you, and hence you can never feel
      affection for him\,/\,her; you always see the other as being in
      opposition. That is why in families people live together for years and
      still don't feel related to each other --- the basic trust and respect
      are missing.
\item When affection is there, I have come to realise that I am related to you
      and you are related to me. There is no doubt on intention, hence no
      irritation, anger or feeling of opposition --- and therefore no fights.
      This is harmony in the family.
\item Affection also \textbf{removes competition}. Competition results when
      there is a lack of affection. When there is affection, I \emph{help the
      other grow}; when I miss this feeling, I try to beat the other and act
      as an opponent. What is naturally acceptable to us is excellence, not
      competition --- and we humans can grow only in relationships.
\item Affection is the starting point of the expansion outward: it starts with
      being related to one, and slowly expands to the feeling of being related
      to everyone, which is \textbf{love} --- and this lays the basis of an
      undivided society.
\end{abul}

\textbf{Role of physical facilities.} Affection is a feeling in `I'; it is
\emph{not} fulfilled by physical facilities. Of the nine values, \textbf{only
Care (mamata) requires physical facilities}, because care is the feeling to
nurture and protect the body of our relative. For affection and the other
feelings, what we need essentially is their \emph{proper understanding}.
Trying to express affection only by giving things is precisely the mistake of
reducing relationship to the exchange of physical facilities.

% ------------------------------------------------------------------------ 39
\Q{39}{Explain the feelings of care, guidance, reverence, glory and gratitude.
  \emph{(Also set as:} Illustrate the natural acceptance values in
  relationship with an example for each --- (i) Reverence (ii)
  Glory.\emph{)}}{\tg{CIE-II 2025}\mmk{5}}

\ans
\begin{abul}
\item \textbf{Care (Mamata)} --- the feeling to \textbf{nurture and protect
      the body of our relative}. We understand a human being as a co-existence
      of the Self (`I') and the Body, and the Body is an instrument of `I'.
      Based on this understanding, we take the responsibility of nurturing and
      protecting the body of our relative. \emph{(This is the only value that
      requires physical facilities.)}
\item \textbf{Guidance (Vatsalya)} --- the feeling of \textbf{ensuring right
      understanding and right feelings in the other} (my relative). We
      understand the need of our Self for right understanding and feelings; we
      also understand that the other is similar to me in the faculty of
      natural acceptance, in the desire for continuous happiness, and in the
      potential of desire, thought and expectation. Right understanding is
      also a need of the other, and the other is related to me; hence I have
      the responsibility to help the other.
\item \textbf{Reverence (Shraddha)} --- the feeling of \textbf{acceptance of
      excellence in the other}. We aspire for continuous happiness, and to
      realise it we must understand harmony at all levels of living and live
      accordingly. When we see that the other has achieved this excellence, we
      have a feeling of reverence for him\,/\,her. \emph{Example:} a teacher
      who has actually understood harmony at all four levels and lives
      accordingly --- I accept that excellence in them, and that acceptance is
      reverence.
\item \textbf{Glory (Gaurava)} --- the feeling for someone who has \textbf{made
      efforts for excellence}. There have been people in history, or even
      around us, who invest their time, energy and belongings to achieve
      excellence and to make others excellent. This gives us a feeling of
      glory for them. \emph{Example:} a person who gave up a comfortable
      career to work for the right understanding of others --- I accept their
      \emph{effort} towards excellence, and that acceptance is glory.
\item \textbf{Gratitude (Kritagyata)} --- the feeling of \textbf{acceptance
      for those who have made effort for my excellence}. Today, what we
      generally call gratitude is a feeling coming out of assistance at the
      level of physical facilities; this feeling is short-lived, since the
      physical facility and the sensation from it are also short-lived. But
      gratitude coming out of someone doing something for my \emph{right
      understanding is permanent}, since the happiness we get from right
      understanding is permanent.
\end{abul}

\begin{keybox}
\textbf{\sffamily Reverence, glory and gratitude --- the distinction that is
examined.} Reverence is for the \emph{excellence achieved}; glory is for the
\emph{effort made} towards excellence; gratitude is for the effort made
towards \emph{my} excellence. The object shifts: their achievement, their
effort, their effort for me.
\end{keybox}

% ------------------------------------------------------------------------ 40
\Q{40}{Define love. How can you say that love is the complete value?
  \emph{(Also set as:} ``The feeling of love lays down the basis of an
  undivided society.'' Explain.\emph{)}}{\tg{SEE}\tg{CIE 2 QB}\mmk{5}}

\ans \textbf{The feeling of being related to all is love (prema).}

\textbf{Why it is the complete value.} This feeling or value is also called
the \textbf{complete value (Purna mulya)}, since it is the feeling of
relatedness to all human beings. The other values are directed towards one
other person --- affection is the feeling of being related to \emph{the
other}, care is towards \emph{my relative's} body, guidance towards \emph{my
relative}. Love is the same feeling extended to everyone, \textbf{without
exception}. Hence nothing more remains to be added; it completes the set of
values in relationship. It starts with identifying that one is related to
\emph{the other} human being (the feeling of affection) and it slowly expands
to the feeling of being related to \emph{all} human beings.

\begin{center}\small\sffamily\color{slate}
Trust \ $\rightarrow$ \ being related to one (Affection) \ $\rightarrow$ \
being related to everyone (Love)
\end{center}

\textbf{How it lays the basis of an undivided society.} The feeling of love
leads to an \textbf{Undivided Society (Akhanda Samaja)} --- it starts from the
family and slowly expands to the world family. Every human being has a natural
acceptance for relatedness up to the world family in the form of love. The
feeling of being related to every human being leads to our participation in an
undivided society. With the understanding of values in human relationships, we
are able to recognise the connectedness with every individual correctly, and
fulfil it. Conversely, if we do not understand the values in relationship, we
are governed by our petty prejudices and conditionings; we treat people as
high or low on the basis of body, wealth or belief. All this is a source of
injustice and leads to a fragmented society, while our natural acceptance is
for an undivided society and universal human order. Love is precisely the
feeling that removes every basis for such fragmentation, and hence it is the
foundation of the undivided society.

\section{C. Harmony in the Society (Chapter 9)}

% ------------------------------------------------------------------------ 41
\Q{41}{What is the comprehensive human goal? Explain how this is conducive to
  sustainable happiness and prosperity for all. How can these four goals
  create harmony in society?}{\tg{CIE 2}\tg{CIE 2 QB}\tg{SEE 2025}\mmk{5--6}}

\ans In order to facilitate the fulfilment of the basic aspirations of all
human beings in society, the following \textbf{comprehensive human goal} needs
to be understood:

\renewcommand{\arraystretch}{1.3}
\arrayrulecolor{rulegrey}
\begin{longtable}{C{7mm} P{45mm} P{26mm} P{62mm}}
\rowcolor{slate}\hdr{\#} & \hdr{Goal} & \hdr{Level} & \hdr{Meaning} \\ \hline
\endfirsthead
\rowcolor{slate}\hdr{\#} & \hdr{Goal} & \hdr{Level} & \hdr{Meaning} \\ \hline
\endhead
\qn{1} & \textbf{Right Understanding} (Samadhana) & In every individual &
  Right understanding is necessary for all human beings. When one does not
  have it, one remains disturbed and also acts so as to create disharmony with
  other human beings and with the rest of nature. \\ \hline
\qn{2} & \textbf{Prosperity} (Samriddhi) & In every family &
  The family is able to identify its needs and is able to
  produce\,/\,achieve more than its requirements. \\ \hline
\qn{3} & \textbf{Fearlessness (Trust)} (Abhaya) & In society &
  Every member of society feels related to everyone else, and therefore there
  is trust and fearlessness. \\ \hline
\qn{4} & \textbf{Co-existence} (Saha-astitva) & In nature &
  There is a relationship and complementarity among all the entities in
  nature, including human beings. \\ \hline
\end{longtable}

\textbf{How the four are related (and hence create harmony):}

\begin{apts}
\item The harmony in society \textbf{begins from the individual}. We need to
      ensure right understanding in the individual as the foundation of
      harmony in society.
\item With right understanding, the \textbf{need for physical facilities in
      the family can be ascertained}. By assessing our needs correctly and by
      producing more than required, the family can be prosperous.
\item Assurance of right understanding in the individuals and prosperity in
      the families, together with the understanding of human relationships,
      leads to harmony and \textbf{trust (fearlessness) in society}. When
      every individual is able to live harmoniously in relationship, and the
      needs of all the families are ensured, fearlessness (mutual trust) in
      society will naturally follow.
\item When human beings with right understanding interact with nature, it will
      be in consonance with \textbf{co-existence} and will be mutually
      enriching.
\end{apts}

\begin{keybox}
\textbf{1. Right Understanding \ $\rightarrow$ \ 2. Prosperity \ $\rightarrow$
\ 3. Fearlessness (trust) \ $\rightarrow$ \ 4. Co-existence}
\end{keybox}

\textbf{Why it is conducive to sustainable happiness and prosperity for all.}
With little exploration, we find that \textbf{all four are required}; we are
not satisfied with anything less. This is the basic minimum requirement to
ensure sustainable happiness and prosperity, and also the maximum we can think
of --- we cannot cut down any of them, and the moment we leave any one out,
there will be a loss of continuity and the goal cannot be achieved. The goal
is not only comprehensive but also \emph{universal}, i.e. equally applicable to
all human beings and for all times. Because it secures right understanding in
the individual (so happiness does not depend on sensation), prosperity in the
family (so the need is limited and met), fearlessness in society (so no one
lives in the fear of another), and co-existence with nature (so production
does not destroy its own base), the happiness and prosperity so achieved are
\textbf{sustainable} rather than temporary.

% ------------------------------------------------------------------------ 42
\Q{42}{Explain the comprehensive human goal. How does fearlessness follow from
  right understanding and prosperity?}{\tg{CIE 2}\mmk{5}}

\ans The four constituents of the comprehensive human goal are as listed in
\textbf{Q41}: right understanding in every individual, prosperity in every
family, fearlessness (trust) in society, and co-existence in nature.

\textbf{How fearlessness follows.} Fearlessness (\emph{abhaya}) in society
means that every member of society feels related to everyone else, and
therefore there is trust and fearlessness. It follows from the first two goals
in this sequence:

\begin{abul}
\item \textbf{From right understanding:} With right understanding, the
      individual is able to recognise that relationship \emph{is} --- that it
      exists at the level of `I' --- and is able to recognise and fulfil the
      values in that relationship (trust, respect, affection\ldots). Since
      trust means being assured that the other inherently wants my happiness
      and prosperity, an individual with right understanding \textbf{does not
      doubt the intention of the other}, and hence has no reason to fear the
      other. When every individual is able to live harmoniously in
      relationship, mutual trust in society naturally follows.
\item \textbf{From prosperity:} With right understanding, the need for
      physical facilities in the family can be ascertained; by assessing needs
      correctly and producing more than required, the family becomes
      prosperous. A family that feels prosperous has \textbf{no motive to
      exploit, hoard or grab} from another family --- on the contrary, the
      feeling of prosperity leads to sharing and becoming an aid to the other.
      Deprivation, by contrast, leads to exploiting the other. Hence when the
      needs of all the families are ensured, the basis of mutual fear between
      them disappears.
\end{abul}

Thus fearlessness is not something to be manufactured separately; it is the
\textbf{natural consequence} of right understanding in individuals plus
prosperity in families. This is why the text states the relation as the
sequence: \textbf{Right Understanding $\rightarrow$ Prosperity $\rightarrow$
Fearlessness (trust) $\rightarrow$ Co-existence.}

Contrast this with our state today: in place of working for fearlessness, we
are working for \emph{strategic powers}. In the name of defence we are
misusing the valuable resources of nature to make weapons and ammunition. We
are becoming increasingly more fearful of each other, paranoid that the other
human being is out to get me --- so most countries are busy preparing for war,
in the hope that more and more competence for war will lead to peace.

% ------------------------------------------------------------------------ 43
\Q{43}{Critically examine the state of society today in terms of the
  fulfilment of the comprehensive human goal.}{\tg{CIE 2 QB}\tg{SEE 2025}%
  \mmk{5--6}}

\ans
\renewcommand{\arraystretch}{1.3}
\arrayrulecolor{rulegrey}
\begin{longtable}{P{40mm} P{112mm}}
\rowcolor{slate}\hdr{Goal} & \hdr{Where we are today} \\ \hline \endfirsthead
\rowcolor{slate}\hdr{Goal} & \hdr{Where we are today} \\ \hline \endhead
\textbf{Right understanding in individuals} &
  We are not really working for this. We are talking so much about
  \emph{information and skills}, but ignoring the need for right
  understanding, ignoring the need to understand happiness correctly, and
  ignoring the need to understand and be in relationship. The programme to
  ensure the very first step is missing --- and the route to prosperity and
  social harmony is through right understanding only. \\ \hline
\textbf{Prosperity in families} &
  We are not able to see that the need for physical facilities is limited and
  that we can have more than what we need. We are only talking about how to
  \emph{generate more wealth}. Our economy, education and market lure
  individuals to generate as much wealth as possible. We feel proud declaring
  the number of billionaires and trillionaires, without ever evaluating
  whether this is taking us to the state of prosperity. \\ \hline
\textbf{Fearlessness (trust) in society} &
  In place of working for fearlessness, we are working for \emph{strategic
  powers}. In the name of defence we misuse the valuable resources of nature
  to make weapons and ammunition. We are becoming increasingly fearful of each
  other; most countries are busy preparing for war in the hope that more
  competence for war will lead to peace. We have organisations like the United
  Nations, but even there we do not have programmes to ensure mutual trust
  among people. \\ \hline
\textbf{Co-existence with nature} &
  Instead of co-existing, we are busy figuring out better ways to exploit
  nature. We have assumed that the goal of technological development is to get
  victory over nature, to subjugate the entities in nature and to disrupt
  nature's cycles. We have disregarded the obvious truth that nature is our
  basic support system, and that disturbing its balance will result in our own
  destruction. \\ \hline
\end{longtable}

\textbf{Conclusion.} We have missed out on the core things. There is a need to
re-align our focus towards ensuring right understanding and relationship. For
this we have to understand the harmony at all levels of our living, and move
from \emph{living with assumptions} to \emph{living with right
understanding}. Each one of us needs to start this journey within, to be able
to contribute towards the rectification of this dangerous situation.

% ------------------------------------------------------------------------ 44
\Q{44}{What are the five dimensions of human endeavour in society conducive to
  manaviya vyavastha? How are they helpful in achieving the comprehensive
  human goal?}{\tg{CIE 2 QB}\mmk{5}}

\ans These five dimensions broadly cover all the activities that are necessary
and fundamental to the harmonious existence of human society:

\begin{apts}
\item \textbf{Education -- Right Living (Siksha-Sanskara).} Education $=$ to
      understand harmony at all four levels of living. Right Living
      (\emph{sanskara}) $=$ the commitment and preparedness to live in harmony
      at all four levels of living. This dimension ensures `right
      understanding' and `right feelings' in every individual.
\item \textbf{Health -- Self-regulation (Svasthya-Sanyama).} \emph{Sanyama} is
      the feeling of responsibility for nurturing, protecting and rightly
      utilising the body; \emph{Svasthya} is the state of the body when it is
      fit to act according to the needs of the Self and there is harmony among
      the parts of the body. Sanyama is the basis of Svasthya.
\item \textbf{Justice -- Preservation (Nyaya-Suraksha).} Justice refers to
      harmony in the relationship \emph{between human beings} --- its
      recognition, fulfilment and evaluation leading to mutual happiness.
      Preservation refers to harmony in the relationship \emph{between the
      human being and the rest of nature}, leading to mutual prosperity ---
      i.e. \textbf{Enrichment (samvardhana), Protection (sanrakshana)} and
      \textbf{Right Utilization (sadupayoga)} of nature.
\item \textbf{Production -- Work (Utpadana-Karya).} Work is the labour that
      the human does on the rest of nature; Production is the physical
      facility obtained out of work. Two questions matter: \emph{what to
      produce?} (physical facilities for nurturing, protecting and rightly
      utilising the body) and \emph{how to produce?} (through a
      \textbf{cyclic (avartansheel) process}, in harmony with nature, ensuring
      that every unit is enriched).
\item \textbf{Exchange -- Storage (Vinimaya-Kosha).} Exchange is the
      exchanging of produce between members of society \emph{for mutual
      fulfilment, not for the madness of profit}; Storage is the storing of
      produce after the fulfilment of needs, \emph{with a view to right
      utilisation in the future, not hoarding}.
\end{apts}

\textbf{How they achieve the comprehensive human goal:}

\begin{cmptable}{Dimension}{Leads to}
Education -- Right Living (Siksha-Sanskara) &
  \textbf{Right understanding} (samadhana) in the individual \\ \hline
Health -- Self-regulation (Svasthya-Sanyama) &
  \textbf{Prosperity} (samriddhi) --- well-being of the body and
  identification of the need for physical facilities \\ \hline
Justice -- Preservation (Nyaya-Suraksha) &
  \textbf{Fearlessness} (abhaya) and \textbf{Co-existence} (saha-astitva)
  respectively \\ \hline
Production -- Work (Utpadana-Karya) &
  \textbf{Prosperity} and \textbf{Co-existence} --- production is done in
  harmony with nature \\ \hline
Exchange -- Storage (Vinimaya-Kosha) &
  \textbf{Prosperity} and \textbf{Fearlessness} --- when we store and exchange
  for mutual fulfilment and not for exploitation \\ \hline
\end{cmptable}

% ------------------------------------------------------------------------ 45
\Q{45}{Explain how production activities can be enriching to all the orders of
  nature. Give any two examples. \emph{(Also set as:} What is the avartansheel
  process of production?\emph{)}}{\tg{CIE 2 QB}\mmk{5}}

\ans On understanding the harmony at all levels of our living, it becomes
evident that there is an inherent balance in nature. So it is only natural
that any production system we design is \textbf{within the framework present
in nature}. When we look at the way nature is organised, two things become
apparent:

\begin{apts}
\item The systems in nature are \textbf{cyclic}, i.e. they are not open-ended.
\item The systems in nature are \textbf{mutually fulfilling or mutually
      enriching}.
\end{apts}

Thus the way to produce is through a \textbf{cyclical (avartansheel) process,
in harmony with nature}: (a) it has to be cyclic, and (b) it has to ensure
that every unit is enriched.

\textbf{Example 1 --- The seed and the soil (nature's own cycle).} When a seed
is planted in soil and water is added, it grows into a tree and in turn bears
leaves, flowers and fruits. The fruits ripen, leaves mature and fall to the
ground and enrich the soil, forming manure by decaying. Seeds are scattered
from the fruit into the soil and once again form a plant and bear fruit. This
way the soil gets enriched, seeds are multiplied and the tree grows ---
everything is regenerated. The process is both cyclic and enriching.

\textbf{Example 2 --- Extending the cycle: guava and jam.} Guavas grow in
nature; we eat them and finally the residue goes back to the soil through
human excreta. We can extend this process by making jam or jelly out of guava
and eating the jam. This making of jam is production, which is essentially an
extension of the cyclic production process already taking place in nature.

\textbf{Example 3 --- Wood and trees.} The amount of wood one person would
require in a lifetime can be obtained from about four full-grown trees. How
many trees can a person plant in a lifetime? Certainly more than four --- it
can even be ten, twenty or a hundred. So, if aware, a human being can be
enriching for nature in a much more effective manner than an animal can.

\textbf{What we do instead.} While nature's processes are all cyclic
(close-ended), our processes are \textbf{acyclic (open-ended)} and deplete
resources. When we burn coal, it is a non-renewable resource --- we can never
again produce it. The utility of the fossil fuels at the bottom of the earth
is to keep the temperature on the earth's surface in a steady state; by
depleting them we tamper with the earth's ability to maintain its temperature,
which is an irreparable damage. And burning them in enormous quantities
pollutes the atmosphere and poisons the air we breathe. Similarly, on
enriching we largely fare `NO': land that has seen heavy use of chemical
fertilisers becomes unfit for agriculture, and pesticides poison our own
bodies as well as animals and birds.

% ------------------------------------------------------------------------ 46
\Q{46}{Describe the concept of an undivided society and the universal human
  order, and explain how both can help create a world
  family.}{\tg{CIE 2 QB}\tg{SEE}\mmk{5}}

\ans
\begin{abul}
\item \textbf{Undivided Society (Akhanda Samaja)} --- the feeling of being
      related to \textbf{every human being}.
\item \textbf{Universal Human Order (Sarvabhauma Vyavastha)} --- the feeling
      of being related to \textbf{every unit}, including human beings and the
      other entities of nature.
\end{abul}

\textbf{Undivided society.} Harmony in the family is the building block for
harmony in society, and harmony in society leads to an undivided society when
we feel related to each and every human being. Today our feelings for our
society have become very limited and each one of us lives in a very small web
of relationships. Our natural acceptance, however, is for relatedness with
all, and we can very naturally expand into the world family. The feeling of
being related to every human being --- i.e. love, the complete value --- leads
to our participation in an undivided society. With the understanding of values
in human relationships, we are able to recognise the connectedness with every
individual correctly, and fulfil it.

\textbf{Universal human order.} When we understand the values in relationship
with the other units in nature too, we are able to recognise our connectedness
with them and fulfil it. This enables us to participate in the universal human
order. Having understood the comprehensive human goal, we are able to be in
harmony not only with human beings but also with the rest of nature; we see
that we are related to every unit in nature and ensure mutual fulfilment in
that relationship.

\textbf{How both help create a world family.} Working on the five dimensions
of human endeavour in the light of right understanding, we are able to work
for an \textbf{orderly living of human society, whose foundational unit is the
family and whose final destination is the world family}. Thus a number of
family units in the form of a village, and a number of villages, integrate
into larger clusters of human society --- expanding in this sequence finally
to a universal human order on this planet. Living in this order, we are able
to plan for the need of physical facilities, the availability of natural
resources and the role of human beings in ensuring the need at the level of
the planet. We are able to work for the inculcation of universally acceptable
human values through education, plan systems to ensure justice for all human
beings, and make policies for the well-being of all.

\begin{keybox}
Individual $\rightarrow$ Family $\rightarrow$ Society $\rightarrow$
Nature\,/\,Existence \hfill $\equiv$ \hfill Family order $\rightarrow$ World
family order
\end{keybox}

% ------------------------------------------------------------------------ 47
\Q{47}{Right understanding in the individuals is the basis for harmony in the
  family, which is the building block for harmony in the society. Give your
  comments.}{\tg{CIE 2 QB}\tg{Scheme}\mmk{5}}

\ans Right understanding in the individuals is the basis for harmony in the
family, which is the building block for harmony in the society.

\begin{apts}
\item The harmony in the society \textbf{begins from the individual}. We need
      to ensure right understanding in the individual as the foundation of
      harmony in the society.
\item With right understanding, the \textbf{need for physical facilities in
      the family can be ascertained}. By assessing our needs correctly and by
      producing more than required, the family can be prosperous.
\item Assurance of right understanding in the individuals and prosperity in
      the families, and the understanding of human relationships, leads to
      \textbf{harmony and trust (fearlessness) in the society}. When every
      individual is able to live harmoniously in relationship, and the needs
      of all the families are ensured, fearlessness (mutual trust) in society
      will naturally follow.
\item When human beings with right understanding interact with nature, it will
      be in consonance with the \textbf{co-existence} and will be mutually
      enriching.
\end{apts}

We may also understand it in the following sequence:

\begin{keybox}
\textbf{1. Right understanding \ $\rightarrow$ \ 2. Prosperity \ $\rightarrow$
\ 3. Fearlessness (trust) \ $\rightarrow$ \ 4. Co-existence}
\end{keybox}

% ------------------------------------------------------------------------ 48
\Q{48}{What is the role of the value system in family harmony? \emph{(Also set
  as:} Explain the difference between response and reaction.\emph{)}}%
  {\tg{SEE}\mmk{5}}

\ans The values (feelings) in relationship are definite and can be recognised.
Their recognition, fulfilment and evaluation lead to \textbf{mutual
happiness} --- which is exactly what justice in the family means. Living with
these feelings is our innate need; problems arise in relationships because we
are unable to ensure the continuity of these feelings.

The role of the value system in family harmony can be seen from the following:

\begin{abul}
\item Relationship is \textbf{already there}; it does not have to be created.
      It only needs to be \emph{understood and fulfilled accordingly} --- and
      the value system is precisely what allows us to recognise and fulfil it.
\item \textbf{Trust is the foundation value.} If trust --- the \emph{adhara
      mulya} --- is shaken, then the whole relationship is disturbed. The
      problem today is that even in families we doubt each other, with the
      result that we behave like enemies, we try to put the other person down,
      and there is a breakdown of relationship.
\item The basic values to be understood in relationship are \textbf{trust and
      respect}. If we have these, the remaining values (affection, care,
      guidance, reverence, glory, gratitude, love) flow quite naturally, and
      harmony in the family is ensured.
\end{abul}

\textbf{Response vs.\ reaction.} If we look at our living today, it is largely
in the \emph{reaction} mode and not in the \emph{response} mode:

\begin{cmptable}{Reaction}{Response}
Doubt on intention & We are able to see that relationship IS, at the level of
  `I' \\ \hline
Irritation & We feel the relatedness with the other at the level of `I'
  \\ \hline
Getting angry & We don't doubt the intention of the other `I' \\ \hline
Fights & We feel a sense of responsibility to improve our own competence and
  the other's competence; we work for mutual fulfilment \\ \hline
\end{cmptable}

For want of proper understanding of relationships, we keep reacting to the
behaviour of the other person and are at the mercy of the situation. Only when
we recognise the relationships in terms of appropriate values do we respond to
every situation and to every person in the right way.

% ------------------------------------------------------------------------ 49
\Q{49}{How can exchange of physical goods be mutually fulfilling? Evaluate the
  motivation of exchange in today's scenario. Why is storage
  required?}{\mmk{5}}

\ans \textbf{Exchange (vinimaya)} refers to the exchange of physical
facilities between the members of society; \textbf{storage (kosha)} refers to
the storage of physical facilities that is left after fulfilling the needs of
the family.

\begin{abul}
\item \textbf{Exchange} --- exchanging of produce for \emph{mutual
      fulfilment} (with a view to mutual fulfilment, \textbf{not the madness
      of profit}).
\item \textbf{Storage} --- storing of produce after fulfilment of needs (with
      a view to \emph{right utilisation in future}, \textbf{not hoarding}).
\end{abul}

\textbf{How exchange can be mutually fulfilling.} Each family has the capacity
to produce more than what it needs for itself. If a family produces wheat, it
can produce for ten families together; another family can similarly produce
cotton for its whole neighbourhood. Summing up all the needs in society,
families are capable of producing \textbf{more than the need}. And then we can
exchange things --- in the form of commodities themselves, or through currency
wherever required. We are exchanging so that \textbf{all of us are able to
fulfil our needs together}. As soon as we are able to recognise the
relationship with the other human being or with the rest of nature, we cannot
think of exploiting anything.

\textbf{Why storage is required.} When we produce more than required, we
exchange for our current needs and store for future needs. This storage is to
be used when production is not taking place, or when some relative of ours
needs it. It is for the \textbf{proper utilisation of the physical facility in
the future}, not with a view to hoard.

\textbf{Motivation of exchange today.} We have developed efficient ways of
selling and buying, sending or receiving money, and investing it to multiply
faster than nature ever could. Sitting with a laptop we can purchase
commodities across the world and invest capital in distant markets; profits
can multiply overnight and one can enter the list of trillionaires without any
physical work; we can store hoards of currency within a digital map. But with
these rising modes of exchange and storage, the \textbf{exploitation of mankind
and nature has shot up}. The disparities in wealth have increased, and the
\textbf{madness for profit has become the general motivation}. Liquidity of
money has of course provided a smooth mode of exchange, but it has created
more problems than solutions. These problems are the outcome of our
mis-perception in visualising money, which is a \emph{notional} entity, to be
the same as \emph{physical facilities}, which are tangible and are our real
needs. It needs to be remembered that \textbf{money is not a need in itself but
only a mechanism to facilitate the exchange of physical facilities}.

\section{D. Further Questions from the 2025--2026 Papers}

% ------------------------------------------------------------------------ 67
\Q{67}{The complete content of respect is to be able to see that ``the other
  is similar to me and we are complementary''. Define complementary in this
  context.}{\tg{CIE-II 2025}\mmk{5}}

\ans The complete content of respect has \textbf{two halves}, and the question
asks about the second.

\textbf{First half --- ``the other is similar to me''.} On the three
evaluations set out in \textbf{Q35}, at the level of `I' the other has the
same basic aspiration (continuous happiness and prosperity), the same
programme of action (to understand and live in harmony at all four levels) and
the same potential (the activities and powers of the Self are continuous and
the same in both of us). That similarity is what makes right evaluation ---
respect --- possible at all.

\textbf{Second half --- ``we are complementary''.} \emph{Complementary} here
means that the difference which does remain between us is \textbf{not a ground
for superiority or inferiority, but a ground for mutual completion}. We are
the same in aspiration, programme and potential; we differ only in
\textbf{competence} --- that is, in how much of that potential each of us has
so far activated, and in our level of \emph{understanding} (not information).
The text is explicit that this difference manifests as a \textbf{meaningful
responsibility}, not as a criterion for ranking:

\begin{apts}
\item If the other has \textbf{better understanding} than me, I want to
      understand from the other --- the other completes what is lacking in me.
\item If the other has \textbf{less understanding} than me, I accept the
      \textbf{responsibility to improve the understanding of the other} --- I
      complete what is lacking in the other.
\end{apts}

So to be complementary is to be \textbf{an aid to the other and not a
hindrance}: each of us supplies what the other lacks, and neither of us is
diminished by the exchange. It is the opposite of competition, in which the
difference in competence is read as a ranking and the other becomes an
opponent to be beaten. Recall from \textbf{Q38} that affection removes
competition precisely because, once I feel related, I help the other grow
instead of trying to beat the other.

\begin{keybox}
\textbf{Similar} $=$ same at the level of `I' (aspiration, programme,
potential). \quad \textbf{Complementary} $=$ the remaining difference, which
is only in competence and understanding, is a basis for helping each other
rather than for ranking each other.
\end{keybox}

Note the consequence. If I see only the similarity, I may feel related but
have nothing to do about the difference. If I see only the difference, I
differentiate --- and, as \textbf{Q36} shows, differentiation on body,
physical facilities or beliefs leads to discrimination and to a fragmented
society. Respect is complete only when both are seen together: the other is
similar to me, \emph{and} where we differ, we complete each other.

% ------------------------------------------------------------------------ 68
\Q{68}{Discuss how right understanding in relationships contributes to harmony
  in the family. Analyse the role of trust, respect and responsibility in
  maintaining family unity, and provide examples to support your
  answer.}{\tg{SEE 2025}\mmk{7}}

\ans \textbf{How right understanding contributes.} Relationship is
\emph{already there} between the Self (`I') and the other Self (`I'); it does
not have to be created, only \emph{understood and fulfilled}
(\textbf{Q27}). Right understanding is exactly what lets us do this: it
enables us to recognise that the feelings in relationship are definite, to
identify them, and to fulfil them. Their recognition, fulfilment and
evaluation lead to mutual happiness --- which is what justice in the family
means. Where right understanding is missing, we reduce the relationship to the
exchange of physical facilities (\textbf{Q30}) and the feelings in `I' go
unaddressed, however well the bodies are provided for.

\textbf{The role of trust (vishvas).} Trust is the \textbf{foundation value
(adhara mulya)}. To trust is to be assured that the other inherently wants
oneself and the other to be happy and prosperous. With right understanding I
separate \emph{intention} from \emph{competence}: intention is always correct
and the same in all of us; only competence is lacking, and competence can be
improved. If trust is shaken, the whole relationship is disturbed --- we
behave like enemies even inside a family.

\emph{Example.} A family member fails to do something he promised. If I am
assured of his intention and see that only his competence was lacking, I do
not get hurt --- I become a help to him, and both of us are comfortable. If I
doubt his intention, we fight, and both of us are unhappy.

\textbf{The role of respect (sammana).} Respect is \textbf{right
evaluation} --- to evaluate the other as he\,/\,she is, at the level of `I',
where the other is similar to me in aspiration, programme and potential. Over,
under or otherwise evaluation all make us uncomfortable; that discomfort is
what we call disrespect. Within a family, differentiation on the basis of the
body (gender, age), of physical facilities (who earns more) or of beliefs is
precisely what produces the generation gap and the sense of opposition between
members.

\emph{Example.} Evaluating a sister as ``only a girl'', or a grandparent as
``too old to understand'', is evaluation at the level of the body, not of `I'.
Evaluated wrongly, they feel disrespected; evaluated as they are, both they
and I are at ease.

\textbf{The role of responsibility.} Once trust and respect are ensured, the
difference in competence between members becomes a \textbf{responsibility}
rather than a ranking (\textbf{Q67}). At the level of feelings this
responsibility takes definite forms named in the text: \emph{care} (mamata)
--- the responsibility of nurturing and protecting the body of my relative;
and \emph{guidance} (vatsalya) --- the responsibility of ensuring right
understanding and right feelings in my relative. Living in this mode is what
the text calls \emph{response} rather than \emph{reaction} (\textbf{Q48}).

\emph{Example.} A younger sibling is doing badly at studies. In the reaction
mode I get irritated and shout, doubting his intention. In the response mode I
accept the responsibility to improve his competence, and I remain related to
him throughout.

\textbf{How this maintains family unity.} With trust there is no doubt on
intention, hence no irritation, anger or opposition --- and therefore no
fights. With respect there is no differentiation, hence no basis for one
member to feel high or low. With responsibility the difference in competence
is worked on together instead of being fought over. Affection --- the feeling
of being related to the other --- can then arise, and affection removes
competition inside the family, because I help the other grow rather than try
to beat the other. This is harmony in the family, and it is why the family is
the \textbf{natural laboratory} in which we gain the assurance that we can
live rightly with human beings at all --- an assurance that then expands from
the family order to the world family order.

% ------------------------------------------------------------------------ 69
\Q{69}{How do justice and preservation contribute to a harmonious society?
  Analyse the shortcomings of focusing solely on punishment in cases of
  injustice, and explain how true preservation ensures prosperity and
  coexistence with nature.}{\tg{SEE 2025}\mmk{6}}

\ans \textbf{Justice and preservation are one of the five dimensions of human
endeavour} --- \emph{Nyaya-Suraksha} (\textbf{Q44}). They are paired because
they are the same relationship-understanding applied in two directions:

\begin{abul}
\item \textbf{Justice (nyaya)} --- harmony in the relationship \emph{between
      human beings}: recognition of values, their fulfilment, evaluation of
      the fulfilment, and mutual happiness ensured. It leads to
      \textbf{fearlessness (trust) in society}.
\item \textbf{Preservation (suraksha)} --- harmony in the relationship
      \emph{between the human being and the rest of nature}, leading to mutual
      prosperity. It has three aspects: \textbf{Enrichment (samvardhana),
      Protection (sanrakshana)} and \textbf{Right Utilisation (sadupayoga)}.
      It leads to \textbf{co-existence with nature}.
\end{abul}

Between them they secure two of the four constituents of the comprehensive
human goal, which is why a society cannot be harmonious without both.

\textbf{Shortcomings of focusing solely on punishment.} The text's definition
of justice makes the shortcoming plain: justice is complete only when
\textbf{all four elements} are ensured, and the fourth is \textbf{mutual
happiness (ubhay-tripti)}. Punishment cannot deliver it.

\begin{apts}
\item \textbf{Punishment addresses only the fourth element, and only
      negatively.} It may register that a value was violated, but it does not
      produce recognition of the value, nor its fulfilment, in either party.
      Three of the four elements are simply left out.
\item \textbf{It cannot produce mutual fulfilment.} In justice there is mutual
      fulfilment for \emph{both} parties. A judgement produces a winner and a
      loser; neither is left feeling related to the other, so the relationship
      is not restored.
\item \textbf{It is a temporary intervention where a continuous need exists.}
      Justice is needed \emph{every moment and in every relationship} --- with
      the child at home, the old grandparent, the maid, close friends, distant
      relations. Courts act after the event and only in the few cases brought
      to them.
\item \textbf{It leaves the actual cause untouched.} Injustice arises from a
      \emph{lack of right understanding} --- from doubting the intention of
      the other, and from evaluating people on body, wealth or belief.
      Punishment does not supply right understanding; it operates on the body
      and on physical facilities, which is the wrong level entirely.
\item \textbf{Hence the text's own question:} will justice be ensured in the
      family, or in the courts of law? Justice has to be \emph{understood and
      lived} in relationship; the legal apparatus can at best contain the
      damage where it has already failed.
\end{apts}

\textbf{How true preservation ensures prosperity and co-existence.}
Preservation is not the mere absence of harm; it is enrichment, protection and
right utilisation together.

\begin{abul}
\item \textbf{It ensures prosperity} because prosperity is the feeling of
      having \emph{more than required} physical facilities, and for all
      physical facilities we are directly or indirectly dependent on nature.
      The continuity of prosperity can therefore be ensured only if our
      production systems are in harmony with nature. Production must follow a
      \textbf{cyclic (avartansheel)} process that enriches every unit
      (\textbf{Q45}) --- as the seed-and-soil cycle does. A person needs the
      wood of about four full-grown trees in a lifetime and can plant ten,
      twenty or a hundred: identify the limited need, produce more than it,
      and both the human and nature gain.
\item \textbf{It ensures co-existence} because co-existence in nature means
      relationship and complementarity among all the entities, human beings
      included. Recognising that we are related to every unit in nature, we
      cannot think of exploiting any of it.
\item \textbf{The contrast with what we do today} shows the mechanism in
      reverse. Our processes are acyclic and open-ended: we burn fossil fuels
      whose actual utility is to keep the earth's surface temperature steady,
      and we cannot produce them again; heavy chemical fertiliser use makes
      land unfit for agriculture. This destroys the very base on which
      prosperity rests, so the exploitation is self-defeating --- nature is
      our basic support system, and disturbing its balance results in our own
      destruction.
\end{abul}

\textbf{Conclusion.} Justice makes society fearless; preservation makes it
prosperous and keeps it in co-existence with nature. Punishment substitutes
for neither, because it works after the fact, on one side only, and at the
level of the body rather than of understanding.

% ------------------------------------------------------------------------ 70
\Q{70}{Explain the concept of harmony in the family. Why is it important, and
  how can it be achieved through right understanding and mutual
  relationships?}{\tg{SEE 2025}\mmk{7}}

\ans \textbf{The concept.} Harmony in the family is the state in which the
values (feelings) in relationship between the members are \emph{recognised,
fulfilled and rightly evaluated}, resulting in \textbf{mutual happiness}. The
family is the \textbf{basic unit of human interaction}; relationship there is
not created but \emph{is}, between the Self (`I') and the other Self (`I'),
and the feelings in it are definite and identifiable --- the nine values, of
which trust is the foundation value and love the complete value
(\textbf{Q31}).

\textbf{Why it is important.}

\begin{apts}
\item \textbf{It is where justice is actually learnt.} The child gets its
      understanding of justice in the family and carries it out into society.
      If the understanding of justice is ensured in the family, there will be
      justice in all our interactions in the world at large.
\item \textbf{It is the natural laboratory.} By living in relationship at all
      times, we gain the assurance that the other is \emph{an aid to me and
      not a hindrance}. Without that assurance we remain shaky in every
      relationship.
\item \textbf{It is the building block of harmony in society.} Right
      understanding in individuals gives harmony in the family, and the family
      is the building block for harmony in society --- and thence for the
      undivided society and the world family order (\textbf{Q47}).
\item \textbf{The alternative is visible around us.} Where it is missing we
      see mistrust, conflict between generations, insecurity in relationships,
      divorce and neglect of older people --- the family-level consequences
      listed in \textbf{Q17}.
\end{apts}

\textbf{How it is achieved --- through right understanding.}

\begin{abul}
\item Recognise that \textbf{relationship IS} and that it exists at the level
      of `I'; the Body is only the means to express or receive it.
\item Recognise that the \textbf{feelings in relationship are definite} and
      can be named --- trust, respect, affection, care, guidance, reverence,
      glory, gratitude, love --- and learn to fulfil them.
\item Understand that \textbf{only care (mamata) requires physical
      facilities}. For all the other feelings what is needed is their proper
      understanding, not more things. Attempting to express affection only by
      giving things is the mistake of reducing relationship to exchange.
\item Move from the \textbf{reaction mode to the response mode}
      (\textbf{Q48}): from doubt-irritation-anger-fights, to seeing that
      relationship is, feeling related, not doubting intention, and taking
      responsibility for improving competence on both sides.
\end{abul}

\textbf{How it is achieved --- through mutual relationships.}

\begin{abul}
\item \textbf{Ensure trust and respect first.} These two are the basic values;
      once they are there, the remaining seven flow quite naturally. Affection
      cannot arise at all until trust and respect are ensured --- which is why
      people can live together in a family for years and still not feel
      related.
\item \textbf{Do not doubt intention.} Assurance of intention plus recognition
      of lacking competence makes me a help to the other; doubt of intention
      puts us in opposition. No problem arises in a relationship unless the
      intention of the other has been doubted.
\item \textbf{Evaluate at the level of `I', not of the body, facilities or
      beliefs.} That removes the ground for differentiation inside the family.
\item \textbf{Let affection remove competition.} With affection I help the
      other grow instead of trying to beat the other; what is naturally
      acceptable to us is excellence, not competition, and human beings can
      grow only in relationship.
\end{abul}

\textbf{Conclusion.} Harmony in the family is not an emotional accident; it is
the predictable result of right understanding applied to a relationship that
already exists. Understand that relationship is, identify the definite
feelings in it, ensure trust and respect, and respond rather than react ---
and mutual happiness in the family follows. That, extended outwards, is the
route from the family order to the world family order.
